\chapter{振动、共振与相图}

\section{从势阱到谐振子}
非退化稳定平衡附近的光滑势能可展开为$V(x_0+\eta)=V(x_0)+k\eta^2/2+O(\eta^3)$，$k>0$。
只保留最低非平凡项，得到$\ddot\eta+\omega_0^2\eta=0$，$\omega_0=\sqrt{k/m}$。
谐振近似要求所访问的位移范围内高阶项远小于二次项；
“振幅小”应相对于势能变化的特征长度来判断。

若有线性阻尼和正弦驱动力，统一写为
\begin{equation}
 \ddot x+2\gamma\dot x+\omega_0^2x=\frac{F_0}{m}\cos\omega t,
 \qquad \gamma\geq0.
 \label{eq:cm-driven}
\end{equation}
无驱动时特征根$s_\pm=-\gamma\pm\sqrt{\gamma^2-\omega_0^2}$。
欠阻尼$\gamma<\omega_0$时
\begin{equation}
 x(t)=A\e^{-\gamma t}\cos(\omega_d t+\varphi),\qquad
 \omega_d=\sqrt{\omega_0^2-\gamma^2}.
\end{equation}
临界阻尼时$x=(A+Bt)\e^{-\omega_0t}$；过阻尼时为两个衰减指数之和。
振幅包络的衰减率为$\gamma$，机械能的精确变化率则为
$\dot E=-2m\gamma\dot x^2$。弱阻尼时周期平均能量近似按$\e^{-2\gamma t}$衰减。

\section{稳态响应：不要混淆三种频率}
取复数约定$x=\Re(\tilde x\e^{\i\omega t})$，稳态解为
\begin{equation}
 \tilde x=\frac{F_0/m}{\omega_0^2-\omega^2+2\i\gamma\omega},\qquad
 x_{\rm st}=A(\omega)\cos(\omega t-\delta),
\end{equation}
其中
\begin{equation}
 A=\frac{F_0/m}{\sqrt{(\omega_0^2-\omega^2)^2+4\gamma^2\omega^2}},\quad
 \delta=\operatorname{atan2}(2\gamma\omega,\omega_0^2-\omega^2).
\end{equation}
$\gamma>0$时瞬态最终消失。$\gamma=0,\omega=\omega_0$时不存在有界正弦稳态，
特解含$t\sin\omega_0t$，不能把响应公式的发散当成有限稳态振幅。

位移振幅的非零频率峰值位于
\begin{equation}
 \omega_{\rm res}=\sqrt{\omega_0^2-2\gamma^2},\qquad
 0<\gamma<\omega_0/\sqrt2.
\end{equation}
阻尼更大时位移振幅从零频率起单调下降。
这与自由衰减频率$\omega_d$及无阻尼频率$\omega_0$不同。
平均输入功率为
\begin{equation}
 \ave P=\frac{F_0^2\gamma\omega^2/m}
 { (\omega_0^2-\omega^2)^2+4\gamma^2\omega^2 }
 =m\gamma\omega^2A^2,
\end{equation}
对固定$F_0$和常数$\gamma>0$，功率峰值在$\omega=\omega_0$。
在这一黏滞阻尼模型中功率半高全宽为$2\gamma$，品质因数定义为$Q=\omega_0/(2\gamma)$。
在共振稳态中，$Q=2\pi E_{\rm st}/\Delta E_{\rm cycle}$，其中$E_{\rm st}=m\omega_0^2A^2/2$，
$\Delta E_{\rm cycle}$是每周期耗散能量。把此能量比用于自由衰减，或用位移峰宽估计$Q$时，通常需要弱阻尼近似。

\section{单摆：线性近似之外}
长度$l$、质量$m$的理想单摆满足
\begin{equation}
 \ddot\theta+\frac gl\sin\theta=0,\qquad
 E=\frac12ml^2\dot\theta^2+mgl(1-\cos\theta).
\end{equation}
当最大摆角$\theta_0\ll1$时，$\mathcal T_0=2\pi\sqrt{l/g}$。
对$0<\theta_0<\pi$，由能量积分并令$\sin(\theta/2)=\sin(\theta_0/2)\sin\psi$得到
\begin{equation}
 \mathcal T=4\sqrt{\frac lg}\,K(k),\quad
 k=\sin\frac{\theta_0}2,\quad
 K(k)=\int_0^{\pi/2}\frac{\d\psi}{\sqrt{1-k^2\sin^2\psi}}.
\end{equation}
小振幅展开为
\begin{equation}
 \mathcal T=\mathcal T_0\Bigp{1+\frac{\theta_0^2}{16}
 +\frac{11\theta_0^4}{3072}+O(\theta_0^6)}.
\end{equation}
振幅越大周期越长；$\theta_0\to\pi$时周期发散，反映摆在顶部附近停留很久。

\section{相图怎样展示运动}
以$(\theta,p_\theta)$为相空间坐标，$p_\theta=ml^2\dot\theta$。
能量曲线直接给出轨迹：低能量为绕稳定点的闭合曲线，$E>2mgl$时摆连续转动，
$E=2mgl$的分界轨道称为分离曲线。
角度相差$2\pi$表示同一位置，所以单摆相空间是角度周期识别后的柱面；
这里只需用这一事实识别图的左右边界，无须先引入一般流形理论。

\begin{center}
\tikzfigure{pendulum-phase-portrait}

小振荡所在的闭轨道（蓝），分离曲线（红），连续转动轨道（黑虚线）。
\end{center}

若单摆受到黏滞阻力矩$-b\dot\theta$（$b>0$），则$\dot E=-b\dot\theta^2\leq0$。
除趋于不稳定平衡点的特殊初态外，轨道最终趋于稳定平衡点。
保守系统中相邻轨道不会因耗散而汇聚；这一区别将在Liouville定理中精确表达。
一个自治的一自由度保守系统能用能量降维求解，单摆的非线性本身并不等于混沌。
周期驱动的非线性摆可出现复杂动力学，但需要通过Poincaré截面、对初值的敏感性等进一步分析，
不能仅凭波形“不规则”便下结论。

\paragraph{复习与阅读}
区分自由振动、瞬态与稳态，以及位移共振与功率共振；用相图判断线性近似的范围。
参阅Taylor的振动与非线性力学、Morin的振动问题。
